\section{System overview}\label{sec:overview}

A person drags a rectangle, the \emph{lens}, over a chat, a feed or a game. The application keeps a screen capture of a
region that is larger than the lens, so that text next to the lens is already read, erased and translated when a drag or a
scroll brings it into view. Figure~\ref{fig:pipeline} shows the data path (blue) and the scroll path (orange).

\begin{figure*}[t]
\centering
\begin{tikzpicture}[
  font=\scriptsize,
  box/.style={draw,rounded corners=2pt,fill=blue!6,minimum width=0.72in,minimum height=0.44in,text width=0.68in,align=center,inner sep=2pt},
  side/.style={draw,rounded corners=2pt,fill=orange!10,minimum width=0.72in,minimum height=0.44in,text width=0.68in,align=center,inner sep=2pt},
  flow/.style={-{Stealth[length=4pt]},thick},
  sig/.style={-{Stealth[length=4pt]},thick,orange!70!black},
  lab/.style={font=\tiny,inner sep=1pt,fill=white}
]
\node[box] (screen) {Screen\\ pixels};
\node[box,right=0.12in of screen] (cap) {Capture\\ SCStream,\\ 8 fps};
\node[box,right=0.12in of cap] (read) {Read\\ Vision lines,\\ Docs XPC blocks};
\node[box,right=0.12in of read] (keep) {Keep?\\ \code{isLiteral}\\ \code{CodeRegion}};
\node[box,right=0.12in of keep] (tr) {Translate\\ engine,\\ batches of 3};
\node[box,right=0.12in of tr] (erase) {Erase\\ flat fill or\\ MI-GAN};
\node[box,right=0.12in of erase] (set) {Re-set\\ fitted size,\\ colour};
\node[box,right=0.12in of set] (ov) {Overlay\\ panel on\\ the lens};
\foreach \a/\b in {screen/cap,cap/read,read/keep,keep/tr,tr/erase,erase/set,set/ov} \draw[flow] (\a) -- (\b);

\node[side,below=0.62in of cap] (diff) {Frame diff\\ tiles of 16 px,\\ shift $\sigma$};
\node[side,below=0.62in of tr] (pred) {Predictor\\ Eqs.~(\ref{eq:shown})--(\ref{eq:spring})};
\node[side,below=0.62in of ov] (wheel) {Wheel events\\ \code{NSEvent}};
\draw[sig] (cap) -- node[lab,right]{frames} (diff);
\draw[sig] (diff) -- node[lab,above]{confirm $\sigma,t_0,t_1$} (pred);
\draw[sig] (wheel) -- node[lab,above]{$(t_k,d_k,\text{kind})$} (pred);
\draw[sig] (pred.north) -- ++(0,0.2in) -| node[lab,pos=0.2,above]{offset $\Delta(t_0,t_1)$ every display frame} (ov.south);

\begin{scope}[on background layer]
\node[draw,dashed,rounded corners=3pt,fill=black!3,inner sep=0.07in,fit=(read)(set)] (seamless) {};
\end{scope}
\node[lab,anchor=south west,fill=none,inner sep=2pt] at (seamless.north west) {once the page has been still for 0.5\,s (``seamless'')};
\end{tikzpicture}

\caption{The pipeline. A captured frame is read into lines and blocks; blocks that must stay as written are dropped;
the rest are translated, their old glyphs erased, and the translation re-set at the same place. Orange: the scroll path.
Wheel events and the frame difference meet in the scroll predictor, which moves the finished blocks every display
frame.}\label{fig:pipeline}
\end{figure*}

\subsection{Capture and change detection}

The screen is captured with ScreenCaptureKit \citep{sckdoc}. The default rate is 8 frames per second, so the frame period is
$1/8 = 125$\,ms. The stream is configured with a minimum frame interval of one over the rate (capped at 30 frames per second)
and a queue depth of 4. Each frame is also reduced to a gray image. Two gray frames $a,b$ are compared tile by tile: a
$16\times16$ tile counts as changed when at least 2 of its pixels moved by more than 24 luma levels. Counting strong
changes, rather than averaging, keeps one edited glyph from being diluted by the rest of its tile and ignores compression
noise.

The pipeline works in two modes. While content moves, it re-reads only the changed band, shows each translation on a plain box
(``fast'') and follows scrolls by shifting existing blocks instead of reading them again. After the page has been still for
$0.5$\,s (the default \code{settleDelay}), it reads paragraphs, erases the source text and paints the translation in its place
(``seamless''), cross-fading from the plain boxes.

\subsection{Reading: lines and blocks}

Text is read with the Vision framework \citep{visiondoc}. A \emph{line} is what the line recogniser returns, with its
rectangle; a \emph{block} is a paragraph, the unit that is translated and repainted. Blocks come from one of two readers.
\begin{enumerate}
\item The \emph{document reader}, a document request of Vision that returns paragraphs, runs in its own XPC service (Docs) so
  that its model compilation, which can take tens of seconds the first time, never stalls the application. The lens does not
  wait for it: a pass that sends a request waits $S=2$\,s and then reads with the line reader while the request runs on.
  A document answer that is there when the line read ends is used; a later one marks the service ready and the lens reads
  once more.
\item The \emph{line reader} followed by a geometric paragraph grouper. The line reads run in the application's own process at
  this revision. The accurate recogniser has $S=3$\,s, the fast one $S=1$\,s.
\end{enumerate}
When successive slow answers show that the document reader is always late, the service is put to rest for $1,2,4,\dots,30$\,s
before the next request, so a service that is always slower than $S$ is not asked again by every recovery read.

\subsection{Translation}

The lines of a block are translated by an engine the person chooses: Apple's Translation framework \citep{transdoc} with a
language pack, or an engine added in the settings. Lines are sent in batches of at most three new lines in reading order.
By default nothing from the screen leaves the machine; text goes to a network engine only if the person adds one.

\subsection{Latency budget}

Table~\ref{tab:budget} lists the stages and the times that the repository records for them. They were measured on different
machines at different dates and are a budget, not a benchmark of the whole path.

\begin{table*}[t]
\centering
\caption{Stages of the data path, where each runs, and the time the repository records for it. The sources are in
\texttt{sources.tsv}.}\label{tab:budget}
\small
\begin{tabular}{@{}p{1.5in}p{1.75in}p{1.35in}p{1.9in}@{}}
\toprule
Stage & Runs in & Recorded time & Note\\
\midrule
Frame period & ScreenCaptureKit stream & $125$\,ms & default $8$ frames per second\\
Document read, warm & Docs XPC service & $0.33$--$0.48$\,s & patience $S=2$\,s derived from it\\
Line read, accurate, warm & application process & $0.25$--$0.88$\,s & patience $S=3$\,s ($\times3$ margin)\\
Line read, fast & application process & $0.06$--$0.10$\,s & Latin scripts only, $S=1$\,s\\
Translation, Apple low latency & on device & $0.15$\,s median per 3 lines & warm, M1 Pro (engine bench)\\
Inpainting, MI-GAN $512^2$ & GPU, Core ML & $13.3$\,ms median & M4 Pro, one patch (inpaint bench)\\
Settle before the seamless pass & timer & $500$\,ms & default \code{settleDelay}\\
Scroll prediction & main thread, display link & every display frame & no read, no model call\\
\bottomrule
\end{tabular}
\end{table*}

The last row is the point of Section~\ref{sec:scroll}. Every other row is a duration that a scroll cannot wait for: a frame
captured at time $t_0$ describes content that, by the time its translations are ready, has moved. The predictor spends no
read and no model call; it only adds up the events that the frame did not yet contain.
